% TOP_PROBLEM: {"id": 3100069, "title": "Least quadratic nonresidue modulo a prime", "category_id": 1, "category": {"id": 1, "name": "number_theory", "display_name": "Number Theory", "description": "Properties of integers, prime numbers, Diophantine equations.", "slug": "number-theory", "order_index": 1, "created_at": "2026-07-31T15:26:25.670Z"}, "status": "partially_solved"}
% FIRST_POSTED: 2026-09-25
% ENTRY_KIND: statement_only
% !TeX program = lualatex
% Definition and sources only; this file is not an LLM solution attempt.
\documentclass[11pt]{article}
\usepackage[margin=25mm]{geometry}
\usepackage{fontspec}
\setmainfont{Latin Modern Roman}
\usepackage{amsmath,amssymb,mathtools}
\usepackage{unicode-math}
\setmathfont{Latin Modern Math}
\usepackage{xurl,hyperref}
\hypersetup{colorlinks=true,urlcolor=blue}
\setcounter{secnumdepth}{0}
\setlength{\parindent}{0pt}
\setlength{\parskip}{5pt}
\setlength{\emergencystretch}{3em}
\begin{document}
\section*{Least quadratic nonresidue modulo a prime}\label{3100069}
\subsection{Definitions and mathematical statement}
For an odd prime $p$, a quadratic nonresidue is an integer $a$ not divisible by $p$ for which $x^2\equiv a\pmod p$ has no solution. Let $n_p$ be the least positive such integer. Vinogradov's conjecture in the selected form asserts
\[
 \forall{\varepsilon}>0\ \exists C_{\varepsilon}>0\ \forall p\text{ odd prime},
 \qquad n_p\le C_{\varepsilon} p^{\varepsilon}.
\]
The constant depends only on ${\varepsilon}$, not on $p$. Equivalently, $n_p=p^{o(1)}$ as primes tend to infinity. The requirement is unconditional, without assuming a Riemann hypothesis. No computable dependence of $C_{\varepsilon}$ on ${\varepsilon}$ or uniform numerical bound is demanded by this notation.
\par\smallskip\textit{Formulation and concept references:} \href{https://arxiv.org/pdf/2404.08380}{[S1]}.

\subsection{Short English statement}
Is the least quadratic nonresidue modulo a prime smaller, up to a constant, than every fixed positive power of that prime?
\subsection{Sources}
\begin{itemize}
\item[S1]E. Carneiro, M. B. Milinovich, E. Quesada-Herrera and A. P. Ramos, "Fourier optimization, the least quadratic non-residue, and the least prime in an arithmetic progression", arXiv:2404.08380 [math.NT], version dated 13 August 2025.
\url{https://arxiv.org/pdf/2404.08380}.
\textit{Formulation source.}
\end{itemize}
\end{document}
